\documentclass[10pt,a4paper]{amsart}
\usepackage[margin=19mm]{geometry}
\usepackage{amsmath,amssymb,mathtools}
\usepackage[scr=rsfs,cal=euler]{mathalfa}
\usepackage{xfrac}
\usepackage[unicode,hidelinks,bookmarks=false]{hyperref}
\newcommand{\ie}{i.e.\ }
\newcommand{\cf}{cf.\ }
\newcommand{\resp}{respectively\ }
\newcommand{\Subsection}{\subsection}
\providecommand{\nobreakdash}{\nobreak\hbox{-}}
\setlength{\emergencystretch}{2em}
\multlinegap0pt
\usepackage{bm}
\usepackage{bbm}
\usepackage{dsfont}
\usepackage{xspace}
\usepackage{cancel}
\usepackage{mathtools}
\usepackage{amsthm}
\makeatletter
\@ifundefined{theorem}{\newtheorem{theorem}{Theorem}}{}
\@ifundefined{lemma}{\newtheorem{lemma}[theorem]{Lemma}}{}
\@ifundefined{proposition}{\newtheorem{proposition}[theorem]{Proposition}}{}
\makeatother
\title[Optimal Lojasiewicz-Simon Exponents and the Intrinsic Singul]{Optimal Lojasiewicz-Simon Exponents and the Intrinsic Singularities of Analytic Functionals}
\author{Tewodros Amdeberhan, Tai Huy Ha}
\date{Source: arXiv:2609.29918v1 (CC BY 4.0)}
\begin{document}
\maketitle

\noindent\emph{Source and licence.} Optimal Lojasiewicz-Simon Exponents and the Intrinsic Singularities of Analytic Functionals. Version arXiv:2609.29918v1 (CC BY 4.0), distributed under the Creative Commons Attribution 4.0 International licence, as recorded in the arXiv licence metadata for this exact version and shown on the abstract page https://arxiv.org/abs/2609.29918v1. Full rights notice: licenses/arxiv-2609.29918-CC-BY-4.0.txt.\par\smallskip

\noindent\textit{Editorial scope.} Counted unit: the Proposition whose source label is \texttt{prop:morse\_bott\_reduction} in the article (source0.tex lines 284--286, source label \texttt{prop:morse\_bott\_reduction}), delivered together with the article's own complete original proof of it (source0.tex lines 289--333, one \texttt{proof} environment of 45 lines). No mathematics is written, repaired or re-derived: statement and proof are the article's own text line for line. Required context retained uncounted: lem:gradient\_identity (lines 159--163); thm:exact\_reduction (lines 219--222). Every definition, display and label of those lines is retained unchanged. Omitted from this package and not redistributed: 1--158, 164--218, 223--283, 287--288, 334--1427. Nothing inside a retained excerpt is omitted or altered: every retained body is delivered exactly as the article's lines are written, so no omission receipt of any kind accompanies this package. Citations: the retained excerpts carry no citation command at all, so no bibliography entry of the article is transcribed and no reference is redistributed. Reference adaptation: none was needed and none was made; every reference inside the delivered unit resolves inside it, so no unresolved reference or citation is produced. Printed slips kept literal: no printed word, symbol, hypothesis or number of the counted statement or of its proof was altered. Layout and class: the article's own class is \texttt{article} with packages including geometry, fontenc, inputenc, lmodern, enumitem, hyperref, cmbright, amssymb, amsmath, latexsym, amsthm, mathtools, graphicx, bbm; the wrapper is the lane's pinned \texttt{amsart} editorial wrapper at 10pt on A4, which restates the theorem-like environments the retained text needs and adds the article's own macro definitions that the retained text needs (none), with extra wrapper packages (bm, bbm, dsfont, xspace, cancel, mathtools). The delivered numbering is the unit's own. Assets: no retained span contains a figure environment, a graphics-inclusion command, a PDF-page inclusion, a table or a TikZ picture; no image file is copied into this package, the package declares no asset dependency and no image file is redistributed with it. Licence: version arXiv:2609.29918v1 of this article is distributed under the Creative Commons Attribution 4.0 International licence, as recorded in the arXiv licence metadata for this exact version and shown on the abstract page https://arxiv.org/abs/2609.29918v1. A full rights notice is delivered at licenses/arxiv-2609.29918-CC-BY-4.0.txt. Characters: every retained excerpt is pure ASCII, so the delivered engine, XeLaTeX with its default Latin Modern text font, reports no missing character.
\begin{lemma} \label{lem:gradient_identity}
For every $\xi \in V$, we have $D\Gamma(\xi) = \rho\bigl(M(c(\xi))\bigr) \in K^*$.
Equivalently, for every $h \in K$, 
$$D\Gamma(\xi)[h] = \langle M(c(\xi)),h\rangle_{Y,X}.$$
\end{lemma}

\begin{theorem} \label{thm:exact_reduction}
Under the hypotheses above, we have
$$\theta_{\mathrm{LS}}(E,0) = \theta_R(\Gamma,0).$$
\end{theorem}

\begin{proposition} \label{prop:morse_bott_reduction}
Under the hypotheses of Theorem~\ref{thm:exact_reduction}, the functional $E$ is Morse--Bott at $0$ if and only if the reduced energy $\Gamma$ is identically zero in a neighborhood of $0 \in K$. 
\end{proposition}

\begin{proof}
Recall that $K = \ker L$ and $X = K \oplus Z$, and that the Lyapunov--Schmidt parametrization is
$c(\xi) = \xi + \psi(\xi)$ for  $\xi \in V \subseteq K$, with $R M(c(\xi)) = 0$. 
Hence $M(c(\xi)) \in Y_K$. By the gradient identity of Lemma~\ref{lem:gradient_identity},
$D\Gamma(\xi) = \rho(M(c(\xi))) \in K^*$. 
Since $\rho:Y_K\to K^*$ is an isomorphism, $D\Gamma(\xi)=0$ if and only if $M(c(\xi))=0$.

\medskip
Assume first that $\Gamma \equiv 0$ near $0$. Then $D\Gamma(\xi) = 0$ for every sufficiently small $\xi \in K$. By this equivalence,
we know $M(c(\xi)) = 0$. Thus $c(V) \subseteq \operatorname{Crit}E$, where $\operatorname{Crit}E = \{x \in U : M(x) = 0\}$.
Conversely, let $x \in \operatorname{Crit}E$ be sufficiently close to $0$. Since $X=K\oplus Z$, we may write uniquely $x=\xi+z$ with $\xi\in K$ and $z\in Z$.
Since $M(x) = 0$, we have in particular $R M(\xi + z) = 0$. 
By the uniqueness part of the implicit function theorem defining $\psi$, the unique sufficiently small $z \in Z$ satisfying this equation is
$z = \psi(\xi)$. Therefore $x = c(\xi)$. Consequently, $\operatorname{Crit}E = c(V)$ near $0$.

\smallskip
Since $c$ has the form $c(\xi) = \xi + \psi(\xi),  \psi(0) = 0$ and $D\psi(0) = 0$, 
it is a real analytic embedding of $V$ as a finite-dimensional submanifold of $X$, since the projection $P:X=K\oplus Z\to K$ satisfies $P\circ c=\operatorname{id}_V$. Moreover, $T_0\operatorname{Crit}E=Dc(0)(K)=K$ because $D\psi(0)=0$.
So, the critical set is a finite-dimensional real analytic submanifold near $0$ and
$T_0 \operatorname{Crit}E = \ker L$. Thus $E$ is Morse--Bott at $0$.

\medskip
Conversely, assume that $E$ is Morse--Bott at $0$. Let $\mathcal{C} := \operatorname{Crit}E$. 
By the Morse--Bott assumption, after shrinking the neighborhood of $0$, $\mathcal{C}$ is a finite-dimensional $C^2$ submanifold of $X$ satisfying
$T_0 \mathcal{C} = K$. In particular, $\dim \mathcal{C} = \dim K$.

\medskip
Let $P : X = K \oplus Z \longrightarrow K$
be the continuous projection associated with the chosen splitting. Consider its restriction
$P|_{\mathcal{C}} : \mathcal{C} \longrightarrow K$. At the origin, $D(P|_{\mathcal C})_0=P|_{T_0\mathcal C}=P|_K=\operatorname{id}_K$.
Since $\mathcal{C}$ and $K$ have the same finite dimension, the finite-dimensional inverse function theorem implies that, after shrinking $\mathcal{C}$, the map
$P|_{\mathcal{C}} : \mathcal{C} \longrightarrow V'$
is a local diffeomorphism onto a neighborhood $V' \subset K$ of $0$.

\medskip
Thus, for every $\xi \in V'$, there is a unique critical point $x_\xi \in \mathcal{C}$ such that
$P x_\xi = \xi$. Writing $x_\xi=\xi+z_\xi$ with $z_\xi\in Z$ and using $M(x_\xi)=0$, we obtain $RM(\xi+z_\xi)=0$.
By the uniqueness in the implicit function theorem defining the Lyapunov--Schmidt correction,
$z_\xi = \psi(\xi)$. Hence $x_\xi = c(\xi)$. Since $x_\xi$ is critical,
$M(c(\xi)) = 0$ for every $\xi \in V'$.
Applying the gradient identity, we have $D\Gamma(\xi) = \rho(M(c(\xi))) = 0$
on $V'$. Therefore $\Gamma$ is constant on the connected neighborhood $V'$. Since
$\Gamma(0) = E(c(0)) = E(0) = 0$, we conclude that
$\Gamma \equiv 0$ on $V'$. The proof is complete.
\end{proof}

\end{document}
